\chapter{序言}

本书是元逻辑的导论，主要面向计算机科学与哲学专业的学生。“元逻辑”之所以如此称呼，是因为它是一门研究逻辑本身的学科。严格意义上的逻辑关注的是有效推理的准则，其符号化或形式化的版本则使用形式语言（例如命题逻辑和一阶逻辑的语言）来呈现这些准则。元逻辑研究的是这些语言的性质，以及使用这些语言的有效推理准则的性质。它探讨诸如如何为这些形式语言的表达式赋予精确意义、如何论证有效推理准则的正当性、各种推导系统（包括它们的计算性质）具有哪些特性等课题。这些问题本身既重要又有趣，因为所研究的语言和推导系统被应用于许多不同领域——尤其是在数学、哲学、计算机科学和语言学中——但它们也可作为研究一般形式系统的范例。我们在这里研究的逻辑语言并非人们感兴趣的唯一语言。例如，语言学家和哲学家对远比命题逻辑和一阶逻辑复杂的语言感兴趣，而计算机科学家则对完全不同\emph{种类}的语言（如编程语言）感兴趣。我们所讨论的方法——如何为形式语言提供语义、如何证明关于形式语言的结果、如何研究形式语言的性质——同样适用于这些情况。

与任何学科一样，元逻辑既包含一系列结果或事实，也包含一套方法和技术，本书涵盖了这两方面。许多学生在本课程之外并不需要了解我们讨论的所有结果，但他们将需要并运用我们确立这些结果的方法。比如说，L\"owenheim-Skolem（勒文海姆-斯科伦）定理在计算机科学中并不常出现，但我们用来证明它的方法却会。另一方面，我们讨论的许多结果确实与某些争论（例如科学哲学和形而上学中的争论）相关。哲学学生在本课程之外可能不需要能够证明这些结果，但他们确实需要理解这些结果是什么——而只有当你仔细思考了确立这些结果所需的定义和证明，你才真正\emph{理解}了它们。在某种程度上，这就是为什么本书涵盖的结果与方法会被推荐——在某些情况下甚至是要求——计算机科学与哲学专业的学生去学习的原因。

本书内容分为三个部分。\Olref[sfr][][]{part} 讨论集合论。逻辑和元逻辑在历史上与所谓的“数学基础”有着非常紧密的联系。数学基础处理的是诸如整数、有理数、实数、函数、空间等数学对象最终应如何被理解的问题。集合论提供了一种答案（还有其它答案），因此集合论与逻辑长期以来一直并行研究。集合、关系和函数在任何形式的形式化研究中也都无处不在，不仅在数学中，在计算机科学和哲学中一些技术性较强的领域也是如此。当然，为了阐述并证明关于逻辑的语义与证明论以及可计算性基础的结果，拥有一套相应的术语是至关重要的。例如，我们将讨论公式的集合、后承关系与可推导关系、谓词符号的解释（它们实际上是关系）、可计算函数，以及它们之间的各种关系和构造。为这些概念引入简记符号，并透彻思考集合、关系和函数的一般性质，将大有裨益。如果你不习惯数学思维和表述数学证明，那么不妨将第一部分关于集合论的内容看作一个训练场：所有基本定义都会给出，并且我们将使用这些定义给出越来越复杂的证明。请注意，理解这些证明——并且能够自己发现和表述它们——或许比理解结果更重要，尤其是在第一部分。如果数学思维对你来说是崭新的，那么仔细思考其中的例子和习题就非常重要。

但是，在第一部分中，我们将确立一个重要的结果。
这个结果——Cantor（康托尔）定理——依赖于一个在整个科学中都堪称最惊人范例的概念分析，即 Cantor 对无穷的分析。
数个世纪以来，无穷一直困扰着数学家和哲学家。
在 Cantor 之前，没有人知道该如何正确地思考它。
许多人甚至认为思考无穷本身就是一个错误，并相信无穷集合这一概念本身就是不融贯的。
Cantor 将无穷变成了一个我们可以融贯地进行研究的对象，并发展出了一整套关于无穷集合的理论——以及我们用来度量无穷集合大小的无穷数。
他表明，存在着不同层次的无穷。
这套关于“超穷”数的理论既优美又错综复杂，我们不会深入探讨；但我们将能够证明，无穷确实有不同的层次，具体来说，存在着“可数的”和“不可数的”无穷层次。
这个结果有着重要的应用，同时它也确实是一种任何自重的数学家、计算机科学家和哲学家都应该知晓的结论。

在\olref[fol][][]{part}中，我们将转向一阶逻辑。
我们将定义一阶逻辑的语言及其语义，即什么是一阶结构，以及一阶逻辑的语句何时在一个结构中为真。
这将使我们能够做两件重要的事情：（1）我们可以用数学精度定义一个语句何时是另一个语句的\emph{逻辑后承}。（2）我们还可以考察，构成一个一阶结构的关系是如何被在其中为真的语句所描述——刻画——的。
这尤其将我们引向对公理化方法的讨论，其中一阶语言的语句被用来刻画某些种类的结构。
接下来我们将关注证明论，我们将考察 Gerhard Gentzen（根岑）在1930年代定义的相继式演算和自然演绎的原始版本。（你的老师可能只选择讲授其中一种，那么之后提到的“推导”和“可推导性”就指代其所选的那个系统。）
后承的语义概念与可推导性的句法概念，为我们精确把握“一个语句可以从另一些语句推出”这一想法提供了两种完全不同的方式。
可靠性与完备性定理将这两种刻画联系起来。
特别地，我们将证明 G\"odel（哥德尔）的完备性定理，该定理指出：每当一个语句是某些语句的语义后承时，它也能从这些语句中推导出来。
一个等价的表述是：如果一个语句集合是\emph{一致的}——即不可能从中证明出矛盾——那么就存在一个使所有这些语句都为真的结构。

完备性定理的第二种表述或许更令人惊讶。
大约在 G\"odel 证明这个结果时（1929年），德国数学家 David Hilbert（希尔伯特）有一个著名的观点：一致性（即无矛盾性）就是数学存在性所需的全部。
换句话说，只要一个数学家能够融贯地描述一个结构或一类结构，那么他就有权相信这种结构是存在的。
当时，许多人觉得这个想法很荒谬：仅仅因为你能不矛盾地描述一个结构，当然并不意味着这样一个结构真的存在。
但这恰恰是 G\"odel 的完备性定理所断言的。
除了这一似非而是——且无疑在哲学上引人入胜——的方面之外，完备性定理还有两个重要的应用，使我们能够证明关于使给定语句为真的结构之存在性的进一步结果。
这就是紧致性定理和 L\"owenheim-Skolem 定理。

在 \olref[tur][][]{part} 中，我们将逻辑与可计算性联系起来。这同样有其历史渊源：David Hilbert（希尔伯特）曾提出逻辑的一个基本问题，即找到一种机械方法，以判定给定的逻辑语句是否有证明。当然，对于命题逻辑而言，这样的方法是存在的：只需检查所有的真值表；由于真值表只有有限多个，该方法最终会给出一个正确答案。然而对于一阶逻辑，这种直接的方法不再奏效，因为可能的结构有无限多个（而且结构本身也可能是无限的）。多年来，逻辑学家们一直在努力寻找更巧妙的方法。Alonzo Church（邱奇）和 Alan Turing（图灵）最终确立了这样的方法并不存在。为了做到这一点，首先必须对什么是“机械方法”给出一个精确的定义。如果有人提出过一种判定程序，人们大概会认可它是一种有效方法。要证明不存在有效方法，你必须先定义“有效方法”，并基于该定义给出一个不可能性证明。这正是 Turing 所做的：他提出了图灵机（Turing machine）的概念\footnote{Turing 自己当然不这么叫它。}，作为机械过程原则上所能完成之事的数学模型。这是利用数学工具对非形式概念进行\emph{概念分析}的又一个例子；它对于计算机科学的重要性，或许与 Cantor（康托尔）对无穷的分析之于数学的重要性同等重要。我们最后一项主要工作将是证明两个不可能性定理：我们将证明所谓的“停机问题”不能被图灵机解决，并最终证明 Hilbert 的（针对逻辑的）“判定问题”也无法被解决。

本书具有数学性质，因为我们讨论数学定义并以数学方式证明结论。但本书又不属于那种需要广博数学背景知识的数学著作。本书内容不要求代数、三角学或微积分知识。我们还特别致力于不要求读者事先熟悉数学的运作方式：事实上，本书的部分目的正是要\emph{培养}理解和证明我们的结论所需的那种推理与证明技能。本书的组织遵循数学惯例，原因在于：这些惯例之所以被发展出来，是因为清晰和精确尤为重要，因此，例如，判断某个陈述是作为论证的结论被断言，还是作为支持其他结论的理由被提出，抑或是旨在引入新的术语，这是至关重要的。所以我们遵循数学惯例，如果某些段落用于引入新术语或符号，就将它们标注为“定义”；而当我们记录一项结果或发现时，就将其标注为“定理”、“命题”、“引理”或“推论”。除了这些惯例之外，我们将使用在一门基础逻辑课程中可能已经熟悉的逻辑证明方法，并且会广泛使用\emph{归纳法}来证明结论。附录中有两章专门讲解这些证明方法。

\paragraph{给教师的说明}
本书中的材料适合一个学期的形式逻辑进阶课程。我在卡尔加里大学讲授的 Logic~II 中，用 12 周的时间覆盖这些内容，尽管我不会像本书这样详细地讲解所有内容。例如，我通常只讨论自然演绎，而略去关于等词的完备性的详细证明。学生已经修过 Logic~I，这门课通常使用 \emph{\href{https://forallx.openlogicproject.org}{forall x: Calgary}} 作为教材，该书使用与我们相同的自然演绎规则，只是采用 Fitch 格式。

本书的最新 PDF 版本可在 \href{https://slc.openlogicproject.org}{slc.openlogicproject.org} 获取，但版本更新频繁。CC BY 许可授权您自行下载和分发本书。为了确保在您使用本书的整个学期中，所有学生拥有相同的版本，您应该这样做：将您决定使用的 PDF 上传到您的学习管理系统（LMS），而不是仅仅给学生链接。您也可以自由地让书店为您打印 PDF，但有些书店可以直接从亚马逊购买并储备平装书。

一阶逻辑的语法、语义及证明系统，均得到 Graham Leach-Krouse 开发的免费在线逻辑教学软件 \emph{Carnap}（\href{https://carnap.io}{carnap.io}）的支持。该软件允许提交并自动批改练习，如自然演绎与矢列演算的推导、为简单理论给出结构，以及符号化练习等。此外，在 \href{https://turing.openlogicproject.org}{turing.openlogicproject.org} 还有一个图灵机模拟器，可用于演示 \olref[tur][][]{part} 部分的内容。其中的例子已在模拟器中预加载。